% !TeX program = pdflatex
% Universite de N'Zerekore - Analyse L1, 2MTH1 0101
% Chapitre 2 : Calcul differentiel - tronc commun des quatre filieres
\documentclass[11pt,a4paper,oneside]{article}
\usepackage[T1]{fontenc}
\usepackage[utf8]{inputenc}
\usepackage[main=french,provide=*]{babel}
\usepackage{lmodern,microtype}
\usepackage{amsmath,amssymb,amsthm,mathtools}
\usepackage{geometry}
\geometry{left=2.4cm,right=2.3cm,top=2.55cm,bottom=2.45cm,headheight=24pt}
\usepackage{graphicx,booktabs,tabularx,array,enumitem}
\usepackage{tikz}
\usetikzlibrary{calc}
\usepackage[most]{tcolorbox}
\usepackage{fancyhdr,lastpage,titlesec,hyperref,needspace}
\definecolor{Marine}{HTML}{10283F}
\definecolor{Bleu}{HTML}{176D8E}
\definecolor{Sarcelle}{HTML}{118F86}
\definecolor{Or}{HTML}{EDAA47}
\definecolor{GrisTexte}{HTML}{54667A}
\definecolor{BleuPale}{HTML}{EFF5F8}
\definecolor{VertPale}{HTML}{EDF7F3}
\definecolor{OrangePale}{HTML}{FFF6E7}
\hypersetup{colorlinks=true,linkcolor=Bleu,urlcolor=Sarcelle,pdftitle={Analyse L1 - Chapitre 2 - Calcul differentiel},pdfauthor={M. Bakary TRAORE}}
\setlength{\parindent}{0pt}
\setlength{\parskip}{5pt plus1pt minus1pt}
\setlist[itemize]{leftmargin=1.35em,itemsep=2pt,topsep=3pt}
\setlist[enumerate]{leftmargin=1.55em,itemsep=3pt,topsep=3pt}
\renewcommand{\arraystretch}{1.18}
\emergencystretch=2em
\titleformat{\section}{\Large\bfseries\color{Marine}}{\thesection}{.7em}{}[\vspace{1mm}{\color{Sarcelle}\titlerule[.8pt]}]
\titleformat{\subsection}{\large\bfseries\color{Bleu}}{\thesubsection}{.65em}{}
\titlespacing*{\section}{0pt}{19pt}{8pt}
\titlespacing*{\subsection}{0pt}{13pt}{5pt}
\pagestyle{fancy}
\fancyhf{}
\fancyhead[L]{\small\color{Marine}\textbf{UNIVERSITÉ DE N'ZÉRÉKORÉ}}
\fancyhead[R]{\small\color{GrisTexte}Analyse L1 \textbullet\ 2MTH1 0101}
\fancyfoot[L]{\scriptsize\color{GrisTexte}Chapitre 2 \textbullet\ Tronc commun}
\fancyfoot[R]{\scriptsize\color{GrisTexte}\thepage\,/\,\pageref{LastPage}}
\renewcommand{\headrulewidth}{.5pt}
\renewcommand{\footrulewidth}{.3pt}
\newtcbtheorem[number within=section]{definition}{Définition}{enhanced,breakable,colback=BleuPale,colframe=Bleu,boxrule=.65pt,arc=1mm,fonttitle=\bfseries,coltitle=white,colbacktitle=Bleu,attach boxed title to top left={xshift=3mm,yshift*=-2mm},boxed title style={boxrule=0pt,arc=.7mm},before skip=10pt,after skip=9pt}{def}
\newtcbtheorem[number within=section]{theoreme}{Théorème}{enhanced,breakable,colback=VertPale,colframe=Sarcelle,boxrule=.65pt,arc=1mm,fonttitle=\bfseries,coltitle=white,colbacktitle=Sarcelle,attach boxed title to top left={xshift=3mm,yshift*=-2mm},boxed title style={boxrule=0pt,arc=.7mm},before skip=10pt,after skip=9pt}{th}
\newtcbtheorem[number within=section]{exemple}{Exemple résolu}{enhanced,breakable,colback=white,colframe=Bleu!65,boxrule=.6pt,arc=1mm,fonttitle=\bfseries,coltitle=Marine,colbacktitle=BleuPale,attach boxed title to top left={xshift=3mm,yshift*=-2mm},boxed title style={boxrule=.5pt,colframe=Bleu!65,arc=.7mm},before skip=9pt,after skip=9pt}{ex}
\newtcolorbox{attentionbox}{enhanced,breakable,colback=OrangePale,colframe=Or!80!black,boxrule=.6pt,arc=1mm,title={\bfseries Point de vigilance},coltitle=Marine}
\newtcolorbox{applicationbox}{enhanced,breakable,colback=VertPale,colframe=Sarcelle,boxrule=.6pt,arc=1mm,title={\bfseries Application aux quatre filières},coltitle=Marine}
\newtcolorbox{objectifbox}{enhanced,breakable,colback=BleuPale,colframe=Bleu,boxrule=.6pt,arc=1mm}
\newcommand{\R}{\mathbb R}
\newcommand{\dd}{\,\mathrm d}
\newcommand{\e}{\mathrm e}
\begin{document}

% Couverture conforme au modele du polycopie fourni.
\hypersetup{pageanchor=false}\thispagestyle{empty}
\begin{tikzpicture}[remember picture,overlay]
 \fill[Bleu] (current page.south west) rectangle (current page.north east);
 \draw[white,opacity=.09,line width=1pt] ([xshift=2.6cm,yshift=5cm]current page.south west) circle (4.1cm);
 \draw[white,opacity=.09,line width=1pt] ([xshift=16.2cm,yshift=23.5cm]current page.south west) circle (5.4cm);
 \draw[Or,opacity=.8,line width=2pt] ([xshift=3cm,yshift=4.5cm]current page.south west)--([xshift=18cm,yshift=4.5cm]current page.south west);
\end{tikzpicture}
\color{white}\vspace*{3cm}
\begin{center}
{\fontsize{19}{23}\selectfont\bfseries UNIVERSITÉ DE N'ZÉRÉKORÉ\par}\vspace{.5cm}
{\large LICENCE 1 -- TRONC COMMUN \quad / \quad 2MTH1 0101\par}\vspace{1.8cm}
{\color{Or}\fontsize{19}{24}\selectfont\bfseries POLYCOPIÉ DE COURS\par}\vspace{.65cm}
{\fontsize{33}{38}\selectfont\bfseries ANALYSE\par}\vspace{.5cm}
{\fontsize{26}{32}\selectfont\bfseries CHAPITRE 2\par}\vspace{.6cm}
{\fontsize{20}{26}\selectfont Calcul différentiel\par}\vspace{.85cm}
{\large 7 heures : 4 h CM + 2 h TD + 1 h TP\par}\vspace{.4cm}
{\normalsize Météorologie \textbullet\ Hydrologie \textbullet\ Génie de l'environnement\par Gestion des Ressources Naturelles\par}\vspace{.65cm}
{\normalsize Enseignant : M. Bakary TRAORE\par}\vspace{.75cm}
{\color{white!80}\normalsize Année universitaire 2026--2027\par}
\end{center}
\clearpage\color{black}\pagenumbering{arabic}\hypersetup{pageanchor=true}\setcounter{page}{1}

\section*{Présentation du chapitre}
\addcontentsline{toc}{section}{Présentation du chapitre}
\begin{objectifbox}
\textbf{Finalité.} Comprendre et quantifier les variations d'une grandeur en fonction d'une variable, analyser ses extrema et construire une approximation locale fiable.

\textbf{Prérequis.} Fonctions usuelles, composition, limites, continuité et calcul algébrique (chapitre 1).

\textbf{Compétences :}
\begin{enumerate}
\item définir une dérivée et calculer les dérivées successives ;
\item vérifier les hypothèses et appliquer Rolle et le théorème des accroissements finis ;
\item étudier les variations, les extrema et la convexité ;
\item obtenir des développements limités et estimer une erreur ;
\item interpréter une vitesse de variation dans les quatre filières.
\end{enumerate}
\end{objectifbox}
\begin{center}\begin{tabularx}{.95\textwidth}{@{}Xc@{}}\toprule
\textbf{Progression du cours magistral} & \textbf{Temps}\\\midrule
Dérivées et règles de calcul & 1 h 20\\
Théorèmes, variations et extrema & 1 h 20\\
Développements limités et applications & 1 h 20\\\midrule
\textbf{Total CM} & \textbf{4 h}\\\bottomrule\end{tabularx}\end{center}
\textit{Supports associés :} Beamer synthétique (14 diapositives), TD 2 avec corrigé indépendant et TP 1 de calcul différentiel (1 h).

\setcounter{tocdepth}{1}\tableofcontents\clearpage

\section{Dérivée et dérivées successives}
\subsection{Définition et interprétation géométrique}
\begin{definition}{Nombre dérivé en un point}{derivee}
Soit $f:D_f\subset\R\to\R$, et $a$ un point intérieur de $D_f$. On dit que $f$ est \textbf{dérivable en $a$} si la limite finie
\[f'(a)=\lim_{h\to0}\frac{f(a+h)-f(a)}{h}\]
existe. Le nombre $f'(a)$ est le coefficient directeur de la tangente au graphe en $(a,f(a))$ :
\[y=f(a)+f'(a)(x-a).\]
\end{definition}
Cette limite est celle des \emph{taux moyens de variation} lorsque les deux instants (ou abscisses) se rapprochent. Si $f$ est dérivable en $a$, elle est continue en $a$.
\begin{exemple}{À partir de la définition}{carre}
Pour $f(x)=x^2$, $\dfrac{f(a+h)-f(a)}{h}=2a+h\to2a$. Donc $f'(a)=2a$ et la tangente en $a=2$ est $y=4x-4$.
\end{exemple}
\begin{attentionbox}
La réciproque ``continue $\Rightarrow$ dérivable'' est fausse : $f(x)=|x|$ est continue en $0$ mais admet deux pentes unilatérales $-1$ et $+1$. Aux bords d'un domaine, on peut étudier des dérivées unilatérales.
\end{attentionbox}

\subsection{Règles de calcul et dérivée d'une composée}
Si $u,v$ sont dérivables, alors
\[
(u+v)'=u'+v',\qquad (uv)'=u'v+uv',\qquad
\left(\frac uv\right)'=\frac{u'v-uv'}{v^2}\quad(v\neq0).
\]
La \textbf{règle de la chaîne} est $\bigl(g(u(x))\bigr)'=g'(u(x))u'(x)$.
\begin{center}\begin{tabularx}{\textwidth}{@{}lX@{}}\toprule
\textbf{Fonction} & \textbf{Dérivée et condition}\\\midrule
$u^\alpha$ & $\alpha u^{\alpha-1}u'$ lorsque les puissances réelles sont définies et dérivables\\
$\e^u$ & $u'\e^u$\\
$\ln u$ & $u'/u$, pour $u>0$\\
$\sin u,\ \cos u$ & $u'\cos u,\ -u'\sin u$\\
$\sqrt{u}$ & $u'/(2\sqrt{u})$, pour $u>0$\\\bottomrule\end{tabularx}\end{center}
\begin{exemple}{Produit et composée}{produit}
Pour $f(x)=(x^2+1)\e^{-x}$,
\[f'(x)=2x\e^{-x}-(x^2+1)\e^{-x}=-(x-1)^2\e^{-x}.\]
L'exponentielle étant strictement positive, $f'(x)\le0$ sur $\R$.
\end{exemple}

\subsection{Dérivées d'ordre supérieur et approximation linéaire}
Lorsque $f'$ est elle-même dérivable, on note $f''=(f')'$, puis $f^{(3)}$, $f^{(4)}$ et plus généralement $f^{(n)}$. Par récurrence, si $f(x)=\e^{2x}$, on obtient
\[f^{(n)}(x)=2^n\e^{2x}\quad(n\in\mathbb N).\]
Pour une petite variation $h$, la formule $f(a+h)=f(a)+f'(a)h+o(h)$ entraîne l'\textbf{approximation linéaire} $\Delta f\approx f'(a)h$. Si $f(t)$ représente une grandeur mesurée dans une unité $U$ et $t$ un temps en heures, $f'(t)$ s'exprime en $U/\mathrm h$.

\section{Théorèmes du calcul différentiel}
\subsection{Rolle et accroissements finis}
\begin{theoreme}{Rolle}{rolle}
Si $f$ est \textbf{continue} sur $[a,b]$, \textbf{dérivable} sur $]a,b[$ et satisfait $f(a)=f(b)$, alors il existe $c\in]a,b[$ tel que $f'(c)=0$.
\end{theoreme}
\begin{theoreme}{Théorème des accroissements finis (TAF)}{taf}
Si $f$ est continue sur $[a,b]$ et dérivable sur $]a,b[$, alors il existe $c\in]a,b[$ tel que
\[ f'(c)=\frac{f(b)-f(a)}{b-a}. \]
\end{theoreme}
\textbf{Pourquoi le TAF est utile.} Si $m\le f'(x)\le M$ pour tout $x\in]a,b[$, alors
\[m(b-a)\le f(b)-f(a)\le M(b-a).\]
En particulier, si $f'=0$ sur un intervalle, $f$ y est constante ; si $f'>0$ partout, $f$ y est strictement croissante.
\begin{exemple}{Vérification des hypothèses}{tafex}
$f(x)=x^2$ est continue sur $[1,3]$ et dérivable sur $]1,3[$. La pente moyenne vaut $\frac{9-1}{3-1}=4$. Comme $f'(c)=2c$, le TAF donne $c=2$.
\end{exemple}
\begin{attentionbox}
On n'applique ni Rolle ni le TAF sans contrôler la continuité sur l'intervalle fermé et la dérivabilité sur son intérieur. Par exemple, $|x|$ sur $[-1,1]$ ne satisfait pas l'hypothèse de dérivabilité en $0$.
\end{attentionbox}

\subsection{Variations, extrema et points critiques}
\begin{theoreme}{Signe de $f'$ et monotonie}{mono}
Si $f$ est dérivable sur un intervalle $I$, $f'\ge0$ sur $I$ implique que $f$ est croissante ; $f'\le0$ implique qu'elle est décroissante. Les inégalités strictes partout donnent une stricte monotonie.
\end{theoreme}
Pour rechercher les extrema : (1) déterminer le domaine ; (2) calculer $f'$ ; (3) repérer les zéros de $f'$ et les points non dérivables ; (4) dresser le tableau de signes ; (5) comparer les valeurs aux points critiques et, sur un segment, aux \textbf{bornes}.
Un passage du signe $+$ au signe $-$ pour $f'$ signale un maximum local ; le passage $-$ au signe $+$ signale un minimum local. Un zéro de $f'$ \emph{ne suffit pas} à établir un extremum : $f(x)=x^3$ possède $f'(0)=0$ mais aucun extremum en $0$.
\begin{exemple}{Variations de $x^3-3x$}{variation}
Soit $f(x)=x^3-3x$. Sa dérivée $f'(x)=3(x-1)(x+1)$ est positive pour $x<-1$ et $x>1$, négative pour $-1<x<1$.
Ainsi $f$ croît sur $]-\infty,-1]$, décroît sur $[-1,1]$, puis croît sur $[1,+\infty[$. Le maximum local vaut $f(-1)=2$, le minimum local $f(1)=-2$.
Sur $[-2,2]$, $f(-2)=-2$ et $f(2)=2$ : le maximum global $2$ est atteint en $-1$ et $2$, le minimum global $-2$ en $-2$ et $1$.
\end{exemple}

\subsection{Dérivée seconde, concavité et convexité}
Si $f''(x)\ge0$ sur $I$, $f$ est \textbf{convexe} ; si $f''(x)\le0$, elle est \textbf{concave}. Si $f'(a)=0$ et $f''(a)>0$ (respectivement $f''(a)<0$), alors $a$ est un minimum (respectivement maximum) local strict.
Un \textbf{point d'inflexion} est un point où la convexité change de sens. Pour $x^3-3x$, $f''(x)=6x$, qui change de signe en $0$ : $(0,0)$ est donc un point d'inflexion. L'égalité $f''(a)=0$ seule ne permet pas de conclure.

\section{Développements limités et approximations locales}
\subsection{Taylor--Young et erreur}
\begin{theoreme}{Formule de Taylor--Young à l'ordre 2}{taylor}
Si $f$ est deux fois dérivable dans un voisinage de $a$ et $f''$ est continue en $a$ (hypothèse suffisante), alors, lorsque $h\to0$ :
\[ f(a+h)=f(a)+f'(a)h+\frac{f''(a)}2h^2+o(h^2). \]
Le polynôme $P_2(h)=f(a)+f'(a)h+\frac12f''(a)h^2$ donne l'approximation locale d'ordre 2.
\end{theoreme}
Le symbole $o(h^2)$ signifie que l'erreur $r(h)$ vérifie $r(h)/h^2\to0$ : l'approximation s'améliore lorsque $h$ se rapproche de zéro.
\begin{theoreme}{Majoration du reste de Taylor--Lagrange}{reste}
Si $f$ est trois fois dérivable sur le segment reliant $a$ et $a+h$ et si $|f^{(3)}(x)|\le M$ sur ce segment, alors
\[ |f(a+h)-P_2(h)|\le\frac{M}{6}|h|^3. \]
\end{theoreme}
\begin{exemple}{Calcul et majoration d'une erreur}{erreur}
Pour $f(x)=x^3-3x$ au voisinage de $a=1$,
\[f(1+h)=-2+3h^2+h^3.\]
À $h=0.2$, l'approximation $-2+3(0.2)^2=-1.88$ diffère de la valeur exacte $-1.872$ de $0.008$. Comme $f^{(3)}=6$, la majoration théorique vaut $6(0.2)^3/6=0.008$.
\end{exemple}

\subsection{Développements limités usuels en zéro}
\begin{center}\begin{tabularx}{\textwidth}{@{}lX@{}}\toprule
\textbf{Fonction} & \textbf{DL à l'origine}\\\midrule
$\e^x$ & $1+x+\frac{x^2}{2}+o(x^2)$\\[1mm]
$\ln(1+x)$ & $x-\frac{x^2}{2}+o(x^2)$ pour $x>-1$\\[1mm]
$\sin x$ & $x-\frac{x^3}{6}+o(x^3)$\\[1mm]
$\cos x$ & $1-\frac{x^2}{2}+o(x^2)$\\[1mm]
$(1+x)^\alpha$ & $1+\alpha x+\frac{\alpha(\alpha-1)}2x^2+o(x^2)$, pour $x$ proche de $0$ et $1+x>0$\\\bottomrule
\end{tabularx}\end{center}
\begin{exemple}{Calculer une limite indéterminée}{limite}
\[\lim_{x\to0}\frac{\e^x-1-x}{x^2}=
\lim_{x\to0}\frac{\frac{x^2}{2}+o(x^2)}{x^2}=\frac12.\]
\end{exemple}
\begin{attentionbox}
Pour un DL au voisinage de $a$, introduire $h=x-a$. Ne pas utiliser un DL de degré 2 pour une erreur d'ordre 3 sans hypothèses sur les dérivées nécessaires.
\end{attentionbox}

\section{Applications transversales aux quatre filières}
\begin{applicationbox}
\textbf{Météorologie.} $T(t)=25+3\sin(\pi t/12)$ ($t$ en heures, $T$ en degrés Celsius) : $T'(t)=\frac\pi4\cos(\pi t/12)$ en $^\circ$C/h. À $t=0$, la température augmente au rythme de $\pi/4$ $^\circ$C/h.

\textbf{Hydrologie.} $Q(t)=5+2t\e^{-t}$ ($t$ en heures, $Q$ en $\mathrm{m^3/s}$) : $Q'(t)=2(1-t)\e^{-t}$. Le débit atteint son maximum pour $t=1$ heure, car $Q'$ y passe de positif à négatif.

\textbf{Génie de l'environnement.} $C(t)=10\e^{-0.2t}$, concentration en mg/L, avec $t$ en heures : $C'(t)=-0.2C(t)$. Une approximation à petite durée $h$ est $C(t+h)\approx C(t)(1-0.2h)$.

\textbf{Gestion des Ressources Naturelles.} $B(t)=100-20\e^{-0.1t}$, biomasse en tonnes, avec $t$ en années : $B'(t)=2\e^{-0.1t}>0$. Le stock augmente, mais sa vitesse de croissance diminue, car $B''(t)=-0.2\e^{-0.1t}<0$.
\end{applicationbox}
\textit{Précision :} ce sont des modèles pédagogiques simples ; l'ajustement de paramètres à des données réelles exigerait une démarche statistique distincte.

\section{Fiche de synthèse : l'essentiel à retenir}
\begin{objectifbox}
\begin{enumerate}
\item \textbf{Dérivée :} $f'(a)=\lim_{h\to0}(f(a+h)-f(a))/h$ ; dérivabilité implique continuité.
\item \textbf{Calcul :} maîtriser somme, produit, quotient et chaîne ; calculer $f''$ si besoin.
\item \textbf{TAF :} sous continuité sur $[a,b]$ et dérivabilité sur $]a,b[$, $f'(c)=(f(b)-f(a))/(b-a)$.
\item \textbf{Variations :} signe de $f'$ ; \textbf{convexité :} signe de $f''$.
\item \textbf{Extrema :} comparer points critiques et bornes sur un intervalle fermé.
\item \textbf{Taylor :} $f(a+h)=f(a)+f'(a)h+f''(a)h^2/2+o(h^2)$ ; vérifier le caractère local et les conditions d'erreur.
\item \textbf{Applications :} écrire les unités et interpréter le signe de la dérivée.
\end{enumerate}
\end{objectifbox}
\textbf{Auto-vérification.} Êtes-vous capable de dériver $\ln(1+x^2)$, d'étudier $x^3-3x$ et de justifier $\lim_{x\to0}(\ln(1+x)-x)/x^2=-1/2$ ?

\section*{Repères bibliographiques}
\begin{itemize}
\item J.-M. Monier, \emph{Analyse MPSI}, Dunod : dérivation, accroissements finis et développements limités.
\item T. M. Apostol, \emph{Calculus}, vol. 1 : concepts de dérivée, variations et applications.
\item E. Kreyszig, \emph{Advanced Engineering Mathematics} : applications physiques et ingénierie.
\end{itemize}
\end{document}
